% ---------------------------------------------------------------------------
% Draft v0.2 (2026-08-20). Narrative draft: prose and real numbers, figures and
% some baseline rows pending. Markers: [PENDING] = the experiment is queued
% or running; [TO RUN] = a planned baseline with no numbers yet;
% [FIGURE] = a figure to be prepared; [TODO] = a detail to verify or
% fill in. Strip this block, the markers, and Appendix B before submission.
%
% The markdown source of record is ../draft.md, frozen at v0.2. From here on,
% edits happen in this file.
% ---------------------------------------------------------------------------

\documentclass[journal]{IEEEtran}

\usepackage{amsmath,amssymb,amsfonts}
\usepackage{authblk}
\usepackage{graphicx}
\usepackage{booktabs}
\usepackage{array}
\usepackage{microtype}
\usepackage{xcolor}
\usepackage[ruled,linesnumbered]{algorithm2e}
\usepackage[hidelinks]{hyperref}

\graphicspath{{figures/}}

\newcommand{\revps}{\ensuremath{\,\mathrm{rev/s}}}
\newcommand{\fsenv}{\ensuremath{f_{\mathrm{e}}}}
\newcommand{\snr}{\ensuremath{\mathrm{SNR}}}
\newcommand{\Jobj}{\ensuremath{J}}

% ---------------------------------------------------------------------------
% DRAFT MARKERS.
% \pending{...} flags a number or claim whose measurement has not yet landed.
% Every such site states what will be measured.
% \wip{...} flags a passage of text which is structurally in place but whose
% content depends on pending experiments and will be rewritten around the
% measured numbers.
% Remove both macros (and every use) before submission.
% ---------------------------------------------------------------------------
\newcommand{\pending}[1]{{\color{gray}\textsf{[pending:~#1]}}}
\newcommand{\wip}[1]{{\color{gray!70!black}\textit{[WIP:~#1]}}}

% Placeholder for a figure whose graphic is not yet generated.
\newcommand{\figplaceholder}[2]{%
  \framebox[\columnwidth]{\rule{0pt}{#1}\parbox{0.9\columnwidth}{\centering
      \pending{#2}}}}

% Placeholder box for a figure whose content is described by a \wip marker in
% the caption.
\newcommand{\figbox}[1]{\framebox[\columnwidth]{\rule{0pt}{#1}}}

\begin{document}

\title{Estimating Rotor Speeds from Drone Ego-Noise with Scarce
  Annotated Data}

\author[1]{Dmitrii~Mukhutdinov \thanks{d.mukhutdinov@qmul.ac.uk}}
\author[1]{Lin~Wang \thanks{lin.wang@qmul.ac.uk}}

\affil[1]{School of Electronic Engineering and Computer Science, Queen
    Mary University of London}

\maketitle

\begin{abstract}
A drone's rotors imprint a harmonic structure on every audio recording made
on board: a rotor spinning at rate $f$ radiates energy at $f, 2f, 3f, \dots$,
and the positions of these harmonics follow the rotor speed exactly.
Knowledge of the instantaneous motor speeds enables informed ego-noise
suppression and acoustic drone monitoring. We study the task of estimating
the four motor speeds of a quadrotor, per time frame, from on-board audio
alone. Annotated data for this task is scarce: the recordings with motor
speed logs available to us amount to about an hour of audio across two drone
models. Neural regression models trained on this data reach low error on
held-out segments of the training recordings and degrade sharply across
microphones, flight regimes, and drone models. A direct probe explains the
degradation: when we scale the frequency axis of the input by 2\%, the
predictions change by 0.03\% on average, so the models estimate speeds from
loudness, spectral shape, and the training prior over speed values.
Augmentations that transform the audio and the speed labels together ---
time-warping and frequency scaling --- raise the models' use of the frequency
cue and improve generalization, at a measurable cost in cruise precision. To
diversify the data further, we build a neural noise generator that is
constrained by construction to produce a harmonic-plus-broadband signal
driven by a given speed trajectory. Adding its output to the training data as
extra recordings makes the models worse, because the synthetic amplitude
patterns offer an easier cue than the real recordings do. We therefore train
on generated data under the same label-transforming augmentation schedule as
on real data, in a generated-first curriculum, which improves on real-only
training for all three architectures we evaluate (by 46\%, 13\%, and 9\% in
validation MSE), with the improvement concentrated in the regime the real
recordings lack most: stopped rotors. The task has no published direct
baselines, so we compare against methods from two adjacent tasks: multi-pitch
tracking (salience-map models, including a variant modified for the required
frequency resolution) and tacholess order tracking (ridge detection,
iterative adaptive Vold--Kalman filtering, and our own two-stage blind
tracking method, which reaches 0.7--1.0 rev/s on cruise segments against
corrected telemetry and 1.0--1.2 rev/s against measured motor speeds on
recordings held out from all tuning). The comparison makes explicit what
neural models trained on scarce data add over signal processing, and what
they lose.
\end{abstract}

\section{Introduction}
\label{sec:intro}

Most drone applications rely on video. Audio recorded on board could support
voice interaction, acoustic scene monitoring, and self-diagnostics, and it is
held back by one dominant obstacle: the ego-noise of the rotors, which buries
other sources at signal-to-noise ratios between 0 and $-30$ dB. The ego-noise
has structure. Each rotor produces a comb of harmonics whose spacing equals
the rotor's rotation rate, plus broadband aerodynamic noise. A system that
knows the four instantaneous motor speeds knows where every harmonic sits,
and can use this knowledge for informed suppression of the noise.

Motor speeds are available from the flight controller on some platforms, and
this channel is often impractical: many commercial drones do not expose the
logs, logged speeds require careful time alignment with the audio, and (as we
show in Section 3.3) the logs themselves carry timing and value errors that
matter at the precision the harmonics demand. This motivates the question of
this paper:

\begin{quote}
Can we train a model that predicts motor speeds from drone audio alone, and
does so robustly --- across microphones, recordings, flight regimes, and
drone models?
\end{quote}

Answering it leads to a study of what neural models learn when the annotated
data is two drones and about an hour of audio. Our contributions:

\begin{enumerate}
\item An evaluation of neural motor-speed regression on two data sources with
  motor speed logs, with a protocol that separates held-out-segment error
  from cross-microphone, cross-regime, and cross-drone generalization
  (Section 5).
\item A probe showing that models trained on this data barely respond to
  frequency scaling of the input, together with an augmentation family
  (time-warping and frequency scaling, both transforming the speed labels
  consistently with the audio) that raises the models' use of the frequency
  cue and improves generalization (Section 6).
\item A structured neural noise generator, constrained by construction to
  harmonic-plus-broadband output driven by a given speed trajectory, with
  two findings we found necessary for its use: the fidelity of its
  harmonics is limited by the accuracy of the motor speed annotations of
  its training data, and aggregate spectral losses select checkpoints with
  degraded harmonic structure, so checkpoint selection needs a per-harmonic
  measurement (Section 7).
\item An evaluation of generated data as training material: added naively, it
  makes models worse; combined with shortcut-removing measures it
  transfers; combined with the label-transforming augmentation schedule in a
  generated-first curriculum, it improves on real-only training for every
  architecture evaluated, with the gain concentrated at stopped rotors
  (Section 8).
\item Baselines from two adjacent tasks, including our own two-stage blind
  tracking method --- a comb-matched search over the full rate range followed
  by refinement against the acoustic comb --- which reaches 0.7--1.0 rev/s on
  cruise segments without any learning and serves as the reference point
  for the neural results (Section 4).
\end{enumerate}

\section{Related work}
\label{sec:related}

\textbf{Drone ego-noise.} \pending{cite: DREGON and its localization line;
ego-noise reduction with microphone arrays; single-channel enhancement
studies; our earlier technical report on drone noise modeling --- replace
with the citable version.} Prior work treats motor speed logs as given side
information for suppression. Estimating the speeds from audio has received
little direct attention.

\textbf{Multi-pitch tracking.} Estimating several concurrent fundamental
frequencies is a mature task in music information retrieval, with
salience-map models as the standard family \pending{cite: deep salience /
multi-f0 CNN models; Basic Pitch}. A quadrotor at cruise is a hard instance
of this task: four fundamentals inside one octave (roughly 55--110 Hz),
frequently within 1 Hz of each other, under broadband noise.

\textbf{Order tracking.} Rotating-machinery analysis estimates shaft-speed
profiles from vibration or sound. Tacholess order tracking extracts the
speed from the signal itself, classically by ridge detection on a
time-frequency representation, and more recently by adaptive filtering
methods such as iterative adaptive Vold--Kalman filtering (IAVKF)
\pending{cite}. The drone case adds two complications: four simultaneous,
near-equal shafts, and flight-regime changes (idle, ramps, cruise, landing).

\section{Task and data}
\label{sec:task}

\subsection{Task definition}
\label{sec:task-def}

Given a mono or multichannel audio clip from a drone-mounted microphone,
predict the four motor speeds, in revolutions per second (rev/s), on a
regular time grid (one estimate per STFT frame). Rotor identity is
unobservable from audio: two rotors at nearly equal speeds can swap labels
freely. All losses and metrics therefore follow the permutation-invariant
training (PIT) convention: the predicted quadruple is matched to the
ground-truth quadruple in the best order before the error is computed. We
report PIT-matched mean squared error (PIT-MSE, rev$^2$/s$^2$) for
training-time comparisons and PIT-matched MAE or RMSE (rev/s) for
interpretable comparisons. One rev/s of error displaces the fortieth
harmonic by 40 Hz, so errors well under 1 rev/s are needed for
harmonic-informed suppression.

Throughout the paper, \emph{flight regime} means a phase of flight --- ground,
idle/warm-up, take-off ramp, cruise, landing --- and we tag segments by the
logged speeds (cruise: all four rotors above 45 rev/s; ground: all rotors
stopped; the remaining powered segments are ramps and warm-up). The \emph{full
flight envelope} means all powered regimes together, from the first motor
spin-up to the last spin-down.

\subsection{Datasets}
\label{sec:datasets}

Two data sources provide drone audio with motor speed logs.

\begin{itemize}
\item \textbf{DREGON} \pending{cite} (public): an 8-microphone array under a
  MikroKopter quadrotor, indoor free flight, six flight recordings, roughly
  half an hour of audio counted per channel. All flight recordings log
  commanded motor speeds; a subset of recordings additionally logs measured
  motor speeds from the motor controllers, and these recordings serve as our
  only tachometer-grade reference.
\item \textbf{MD2} (unpublished; recorded by Michael Clayton at Queen Mary
  University of London, first described in \pending{our technical report ---
  replace with the citable version}): two 8-channel outdoor free-flight
  recordings of a DJI Matrice 100 with an on-top circular microphone ring,
  $\sim$5 minutes of 8-channel audio ($\sim$40 minutes counted per channel),
  with DJI SDK telemetry logs.
\end{itemize}

For training we mix these recordings with LibriSpeech utterances at SNRs
from $-30$ to 0 dB, resampled to 16 kHz, generating mixtures on the fly so
every epoch sees fresh combinations. Between them, the two sources contain a
few dozen distinct speed trajectories. This is the scarcity that drives the
rest of the paper.

\subsection{Splits and training regimes}
\label{sec:splits}

Every result in this paper uses one fixed three-way split of the real
recordings (Table~\ref{tab:splits}). The validation split is frozen: 37
non-overlapping 8-second clips spanning the full flight envelope of both
drones, and every number we report — for our models and for every baseline —
comes from the same clips under the same per-frame permutation-invariant
matching. The test split is reserved: two further DREGON recordings from the
same room as the validation clips, and two further MD2 flights whose
telemetry we aligned and calibrated with the procedure of
Section~\ref{sec:annotation}. No experiment in this paper reads the test
split; it exists so that the final model selection can be scored once,
after all development decisions are frozen. \pending{final test numbers}

\begin{table}[t]
\centering
\small
\begin{tabular}{lll}
\toprule
Split & DREGON & MD2 \\
\midrule
Train & 5 flight recordings (room 2) & 1 flight (FLY125) \\
Validation (frozen) & 3 recordings (room 1) & 1 flight (FLY124) \\
Test (reserved) & 2 recordings (room 1) & 2 flights (FLY103/108) \\
\bottomrule
\end{tabular}
\caption{The three-way split of the real recordings. The two test-side
DREGON recordings share the room and rig with the validation clips, so
test performance measures generalization to unseen recordings, and
cross-room generalization is measured between train and validation.}
\label{tab:splits}
\end{table}

Training data passes through one of five named regimes, and every model in
this paper is tagged with the regime it was trained under:

\begin{itemize}
\item \textbf{R1 — architecture search.} Online mixtures of the training
  recordings with LibriSpeech speech; light augmentation (random gain,
  polarity, channel drop at $p{=}0.5$ after a warm-up). The regime of the
  architecture search in Section~\ref{sec:overfit}; its validation split
  predates the frozen one above, so its numbers rank architectures and are
  not comparable with later tables.
\item \textbf{R2 — final real-only.} The full flight envelope of the
  training recordings, frequency-scaling augmentation on every chunk
  (the comb and its labels scaled together by $\alpha \in [0.7, 1.3]$),
  time warping, gain and polarity; plus a zero-labeled silence arm
  ($\sim$17\% of chunks: room tone, colored noise floors up to flight
  level, low-frequency rumble) and a reference floor on the speech
  scaling so that silent chunks still carry speech at a normal level.
  The last two additions exist because the real recordings hold only
  26 seconds of rotor silence, all of it quiet — models trained without
  the additions learn ``quiet input means zero'' and fail on zero-speed
  frames that carry sound.
\item \textbf{R3 — generator+comb curriculum.} Pre-train on synthetic
  noise (the neural generator of Section~\ref{sec:generator} plus an
  analytic harmonic comb, half each), then fine-tune on the real regime.
\item \textbf{R4 — comb-only curriculum.} As R3 with the analytic comb as
  the only synthetic source.
\item \textbf{R5 — mixed single-stage.} One training pool: half real,
  one quarter generated, one quarter comb, with the R3 augmentation
  schedule and no curriculum.
\end{itemize}

\wip{The reported R3--R5 runs used the R2 recipe without the silence arm as
their real component; the final tables re-train them on the full R2 recipe
so that the regimes differ only in their synthetic ingredient.}

\subsection{Annotation quality}
\label{sec:annotation}

The logs are imperfect in ways that matter at harmonic precision. Our
measuring tool for all three corrections below is the \emph{harmonic
reconstruction residual}: fit a harmonic model (a Vold--Kalman decomposition,
Section 4.2) to the audio at candidate speed trajectories, and measure the
residual energy; a candidate that places the comb closer to the true
harmonics leaves a smaller residual.

\begin{itemize}
\item \textbf{Timing.} The MD2 telemetry is late relative to the audio, and
  the error grows linearly with time (a sample-clock dilation). We measured
  and corrected both constants per recording (residual timing error
  $\sim$3--5 ms RMS after correction).
\item \textbf{Values.} The MD2 logged speeds are $\sim$0.7\% low; we apply a
  measured multiplicative correction. For DREGON we measured a bias between
  0.35\% and 0.85\%, depending on the estimation method; no independent
  measurement exists that would decide between the candidate corrections, so
  we leave the DREGON labels unchanged and carry the range as an uncertainty
  floor (0.3--0.7 rev/s at cruise speeds) under every DREGON number in this
  paper.
\item \textbf{Refinement.} On top of the corrected logs we compute refined
  annotations: per-window speed trajectories obtained by gradient-based
  minimization of the harmonic reconstruction residual, starting from the
  telemetry. The procedure's precision floor, estimated on the recordings
  with measured speeds, is $\sim$0.2 rev/s; the same floor applies wherever
  the procedure reappears (Section 4.2, stage 2).
\end{itemize}

Section 7 shows that the difference between raw and refined annotations is
decisive for the noise generator.

\begin{figure}[t]
  \centering
  \figbox{7em}
  \caption{\wip{FIGURE: spectrogram with telemetry vs refined harmonic
      overlays}}
  \label{fig:comb-overlays}
\end{figure}

\section{Baselines from adjacent tasks}
\label{sec:baselines}

The task has no published direct baselines. We assemble baselines from the
two adjacent fields, so that the neural results of Sections 5--8 can be
judged against existing methods. Two evaluation protocols appear in this
paper and their numbers are never comparable across protocols: the
\emph{mixture protocol} (fixed held-out spans of the training recordings,
mixed with speech; Sections 4.1 and 5) and the \emph{window protocol}
(16-second windows of the raw recordings spanning both drones and all flight
regimes, per-window PIT-matched MAE against corrected telemetry;
Sections 4.2 and 6). Each table names its protocol.

\subsection{Multi-pitch tracking}
\label{sec:multipitch}

We adapt two salience-map models to the rotor band and train them with
binary cross-entropy against binary salience targets placed at the
ground-truth fundamental bins: a CNN over a harmonic CQT front-end
\pending{cite: multi-f0 salience CNN} and Basic Pitch \pending{cite}.
Per-frame salience peaks are linked into four trajectories by optimal
bipartite (Hungarian) assignment. Off the shelf, both models fail on this
data, and the failure is informative: on DREGON validation, 55\% of frames
have two rotors within 1 Hz of each other --- comparable to the $\sim$0.9 Hz
bin spacing of the standard salience grids --- so the two peaks merge and the
tracked trajectories collapse toward their mean. A modified variant with the
input band narrowed to one octave (55--110 Hz) and a super-resolution output
head (0.15 Hz/bin, trained end-to-end) resolves the near-unison rotors
(per-rotor MAE spread compresses from 1.2--5.5 Hz to 1.8--2.9 Hz).

\begin{table}[t]
  \centering
  \caption{Mixture protocol, DREGON validation (30 clips $\times$ 8
    channels), PIT-matched. $R^2$ is averaged per clip, so rows imply
    slightly different target variances; RMSE is the comparable column.
    $^{\dagger}$Needs no training; scored on the full-envelope validation
    split of Section 5 (both drones), cruise frames.}
  \label{tab:multipitch}
  \footnotesize
  \setlength{\tabcolsep}{4pt}
  \begin{tabular}{@{}p{0.56\columnwidth}cc@{}}
    \toprule
    Method & RMSE (rev/s) & $R^2$ \\
    \midrule
    Direct regression (best neural model of Section 5, same data)
      & 1.62 & 0.93 \\
    Salience CNN (standard grid) & 6.30 & 0.19 \\
    Salience CNN (narrow band + super-resolution) & 4.03 & 0.57 \\
    Basic Pitch (standard grid) & 23.2 & $-16.2$ \\
    Basic Pitch (narrow band + super-resolution) & 11.7 & $-3.2$ \\
    Inverse harmonic clustering (training-free)$^{\dagger}$
      & 23.3 & --- \\
    \bottomrule
  \end{tabular}
\end{table}

After the resolution fix, the best salience model still trails direct
regression by a factor of $\sim$2.5. Salience models designed for discrete
musical notes are a poor fit for four continuous, crossing, near-unison
fundamentals.

The last row of Table~\ref{tab:multipitch} is a recent classical estimator
that needs no training: inverse harmonic clustering by optimal transport
\pending{cite: Bj\"orkman \& Elvander 2026, arXiv:2508.02471}. The method
fits a nonnegative line spectrum to the sample autocovariance and moves its
mass onto candidate fundamentals with a transport cost that makes a
sub-octave placement expensive. With its published configuration it locks
onto the broadband noise floor and scores 38--45 rev/s. Two changes adapt
it to this signal: an analysis floor of 150 Hz, so that the fundamentals are
inferred from their harmonics instead of from the noisy fundamental band,
and 1-second analysis frames. The adapted estimator scores a cruise
PIT-matched MAE of 16.3 rev/s. The adaptation removes most octave errors:
60\% of reported pitches land within 3\% of a true rotor speed (43\% before)
and 9\% on a sub-octave of one (22\% before). The remaining error is a
coverage error: the estimator finds on average 1.4 of the four distinct
speeds per frame, because its sparsity term makes it cheaper to absorb the
harmonics of a second rotor into an existing pitch column than to open a new
one, and four near-unison rotors are exactly that case. Every learned model
in this paper and every tracking method of Section 4.2 reaches an error
smaller by roughly an order of magnitude, so this row serves as the
training-free floor of the task.

\subsection{Tacholess order tracking}
\label{sec:order-tracking}

Order tracking methods need no training data, which makes them a natural
reference point for a learning approach that has almost none. All numbers
below use the window protocol.

\begin{itemize}
\item \textbf{Ridge detection} on a time-frequency representation, the
  classical tacholess method \wip{TO RUN as a standalone row}.
\item \textbf{IAVKF} \pending{cite}: our reimplementation (the method
  extracts one harmonic component at a time --- ``peeling'' --- with adaptive
  bandwidth) \wip{TO RUN: comparison table}.
\item \textbf{Our two-stage method.} Stage 1 finds four speed trajectories
  blindly: a comb-match score over the full plausible rate range is tracked
  through time by dynamic programming, with explicit handling of ramps and of
  the near-unison rotor pairs. Stage 2 refines each trajectory by minimizing
  the harmonic reconstruction residual of Section 3.3 over the trajectory: a
  coupled Vold--Kalman decomposition at the candidate speeds is fitted, its
  residual differentiated with respect to the trajectory, and the trajectory
  updated until the reconstruction aligns with the acoustic comb to a
  fraction of the comb spacing. The stage-2 precision floor is the
  $\sim$0.2 rev/s of Section 3.3. On DREGON recordings held out from all
  tuning, scored against the measured motor speeds, the full blind method
  reaches 0.97--1.22 rev/s.
\end{itemize}

\begin{table}[t]
  \centering
  \caption{Window protocol, PIT-MAE, rev/s. The DREGON uncertainty floor of
    Section 3.3 --- 0.3--0.7 rev/s at cruise --- sits under both DREGON
    columns: DREGON numbers this small measure agreement with telemetry, and
    absolute accuracy is bounded by the floor. The oracle row is better on
    MD2 and worse on DREGON because the blind search can escape a local
    optimum that the telemetry initialization commits to.}
  \label{tab:order-tracking}
  \footnotesize
  \setlength{\tabcolsep}{4pt}
  \begin{tabular}{@{}p{0.52\columnwidth}cc@{}}
    \toprule
    Method & DREGON cruise & MD2 cruise \\
    \midrule
    Two-stage blind (ours) & 0.69 & 1.03 \\
    Two-stage, initialized from telemetry (oracle) & 0.85 & 0.78 \\
    Best neural rows (Section 6 training setup) & 2.5--2.9 & 2.3--2.4 \\
    \bottomrule
  \end{tabular}
\end{table}

This comparison is the reference for the rest of the paper: on steady
flight, a signal processing method with no training data is 2--4$\times$ more
precise at cruise than the neural models. The value of the neural models
must then come from the other columns of the comparison --- coverage of ramps
and idle, single-pass runtime, robustness where the comb assumption
degrades --- and Section 9 collects those columns.

\section{Neural models overfit the scarce data}
\label{sec:overfit}

\subsection{How the three architectures were selected}
\label{sec:archsearch}

We searched over 26 variants of one model family under regime R1: a shared
residual convolutional encoder with squeeze-and-excitation blocks and a
learned frequency-attention pool, combined with different input features and
temporal heads. The variants cover six groups: the bidirectional-GRU
baseline head; temporal-head swaps (Transformer, windowed attention, wider
GRU, dilated temporal convolutions); input-feature changes (multi-resolution
spectrograms, wavelet branches, magnitude+phase channels); pooling changes;
ten causal/streaming head variants (unidirectional GRUs at several widths,
normalizations and dropouts, fully causal front-ends); and seven
frequency-dilated backbone variants. Every variant trained under two arms —
fixed pre-mixed data, and online mixing — with the same optimizer, loss
(PIT-MSE) and budget, 52 runs in total.

Three findings shaped the selection. First, online mixing is itself an
architectural leveler: 21 of 26 variants improved when mixtures were
generated on the fly, and the attention heads improved dramatically (the
Transformer head moved from diverging on fixed data to ranking among the
best online). Second, causal heads are usable online and unstable offline;
the unidirectional GRU with hidden width 128 became the best single run of
the sweep (PIT-MSE 7.33 on the sweep's validation split). Third, fully
causal front-ends underfit badly and were dropped; head-only causality is
the workable streaming compromise. We kept three architectures — the
convolutional network with the bidirectional GRU head, the Transformer
head, and the causal GRU-128 head — one non-causal compact model, one
attention model, and one streaming-capable model.

A caution attaches to the ranking itself: under later, stricter matched
protocols (a shared stream, front-end and schedule, validated on the frozen
full-envelope split), the ordering among the three winners reshuffles, with
the Transformer variants ahead of the causal GRU. Rankings from a single
regime and split do not transfer; this is why the rest of the paper carries
all three architectures through every experiment instead of a single
champion. \wip{TODO: parameter counts, STFT configuration, and output frame
rate in a small table}

All three take STFT features of 1-second audio chunks and emit four speeds
per frame, trained with PIT-MSE on online-generated mixtures.

Generalization is much worse, in three nested ways:

\begin{itemize}
\item \textbf{Across microphones.} Models trained on one microphone of the
  8-channel arrays degrade on the others; treating channels as extra training
  samples restores parity at no architectural cost.
\item \textbf{Across flight regimes.} Models trained on cruise-filtered data
  fail on ground and ramp segments in proportion to how far the speeds sit
  from the cruise prior: per-regime PIT-MSE reaches 2450 rev$^2$/s$^2$ on
  ground segments, against $\sim$15 at cruise. Keeping the full flight
  envelope of the same real recordings in training, take-off ramp included,
  cuts the all-regime error from 338 to 80 rev$^2$/s$^2$ (Transformer) with
  no new data.
\item \textbf{Across drones.} A model trained on DREGON alone scores 7.96
  rev/s PIT-RMSE on an MD2 recording. Adding the other MD2 recording to
  training drops the error on the held-out one to 1.63 rev/s. One recording
  of a new drone suffices, and we observe no zero-shot transfer between drone
  models.
\end{itemize}

The capacity ordering confirms overfitting as the mechanism: the Transformer
head, largest of the three, has the worst held-out error before augmentation
(11.76 PIT-MSE, against 9.71 for the smallest under the same setup), and it
gains the most from every diversification below.

\section{Augmentation, and what the models actually read}
\label{sec:augmentation}

\textbf{The probe.} We resampled validation audio to scale its frequency axis
by $\times 1.02$ --- which shifts every harmonic, and hence the apparent motor
speeds, by 2\% --- and measured the mean relative change of the PIT-matched
predictions. The ideal response is 2\%. The measured response is 0.03\%,
across architectures. The models' predictions are driven by loudness,
spectral shape, and the training prior over speed values; the harmonic
positions contribute two orders of magnitude less than they should. This one
number explains both the good held-out results (the prior is correct on data
that resembles training) and the sharp cross-condition degradation (the prior
is wrong everywhere else).

\begin{figure}[t]
  \centering
  \figbox{7em}
  \caption{\wip{FIGURE: predicted speed vs input frequency scaling, per
      architecture}}
  \label{fig:freq-probe}
\end{figure}

\textbf{Augmentations that transform the labels.} Standard audio
augmentations (gain, polarity, channel drop) keep the speed labels fixed. Two
augmentations transform the audio and the labels together and thereby
manufacture genuinely new audio-and-speed pairs:

\begin{itemize}
\item \textbf{Time-warping}: resample the noise chunk by a slowly varying
  rate $\alpha(t)$ ($|\alpha - 1| \le 0.12$) and warp the speed labels
  identically. Same spectral content, new valid trajectory. Held-out
  PIT-MSE: Transformer $11.76 \to 8.74$ ($-26\%$), convolutional
  $9.71 \to 8.85$, GRU unchanged.
\item \textbf{Frequency scaling}: resample the chunk by a constant
  $\alpha \in [0.75, 1.3]$ and multiply the labels by $\alpha$. This moves
  the whole comb and the mean of the label distribution together, so a model
  that ignores harmonic positions cannot fit the augmented data. (The range
  is asymmetric so that the scaled speeds stay inside the physically
  plausible band \wip{TODO: verify the stated reason}.)
\end{itemize}

Frequency scaling has a consistent, two-sided effect. For the Transformer
setup, full-envelope PIT-MSE improves from 63.7 to 37.6 while DREGON cruise
MAE worsens from 2.48 to 3.23 rev/s (window protocol); the same trade
appears on every architecture we applied it to. The models can be pushed to
read frequency, and the push costs part of the prior-driven precision at
cruise. Augmentation recombines the recordings we have; new drone timbres,
microphone placements, and speed dynamics have to come from somewhere else.
This motivates synthetic data.

\section{A noise generator constrained to harmonic structure}
\label{sec:generator}

We build a generator with the harmonic structure hard-wired: for each rotor,
an oscillator bank at multiples $k \cdot f_r(t)$ of the given speed
trajectory, up to order $k = 80$, with learned time-varying amplitudes, plus
a learned broadband component; both are propagated to the microphone
positions of the target array by free-field delay and distance attenuation
with learned per-microphone corrections. A per-drone embedding (optionally
with per-rotor sub-embeddings) conditions the learned parts. By construction
the generator can only place harmonic energy on the commanded comb, and it
produces unlimited multichannel noise with exact speed labels.

We evaluate the generator's harmonic fidelity with a \emph{per-harmonic line
measurement}: on validation audio, the floor-subtracted power of the
generated signal at each harmonic of the reference trajectory, compared with
the same quantity measured on the real recording. Aggregate spectral
distances cannot see this quantity --- harmonics are sparse, so a generator
that drops every line above order 50 barely moves a multi-resolution STFT
distance.

Two findings from this instrument:

\textbf{Annotation accuracy limits the harmonics.} Trained on data whose speed
labels carry a constant $-0.54\%$ bias (one of the measured candidate values
of the DREGON bias, Section 3.3), the generator's harmonics above order
$\sim$50 collapse: line power drops by 8.6 dB relative to training on exact
labels. At high orders the biased comb misses the true line positions by more
than a line width, the harmonic bank can no longer explain the observed
energy, and the optimizer reassigns that energy to the broadband component.
Training on the refined annotations of Section 3.3 restores the line
\emph{power} at orders 10--49 on real data. Line \emph{sharpness} above order
$\sim$25 remains limited by the training objective itself: the longest window
of the multi-resolution STFT loss is shorter than the beat period of adjacent
harmonics there, so the loss provides no gradient for line width at those
orders, whatever the label quality.

\textbf{Aggregate losses select the wrong checkpoints.} Across training
epochs, the aggregate spectral distance keeps improving while the line
measurement peaks early and then erodes; the correlation between the two
across epochs is negative ($r = -0.37$ in the largest run; the same pattern
appeared in all four training runs where we tracked both). A generator picked
by the standard validation loss has visibly and audibly washed-out harmonics.
We therefore select checkpoints by the line measurement on validation audio
of both drones. \wip{PENDING: the line-vs-loss checkpoint comparison on both
drones}

\begin{figure}[t]
  \centering
  \figbox{7em}
  \caption{\wip{FIGURE: aggregate loss and line measurement vs epoch;
      spectrograms and audio of the two selections}}
  \label{fig:line-vs-loss}
\end{figure}

\section{Generated data as training material}
\label{sec:generated-data}

Three experiments, in the order we ran them.

\textbf{Added naively, generated data hurts.} Adding generator output to the
real training pool as extra noise recordings increased validation PIT-MSE by
27\% \wip{TODO: recover the absolute pair for this number}. The synthetic
amplitude patterns follow the speed trajectory by construction, which hands
the model the loudness cue of Section 6 in a purer form than real
recordings do.

\textbf{With the shortcut removed, generated-only training transfers.} Two
measures: mix the generator's output in even proportion with an analytic
comb --- a synthetic harmonic signal whose per-harmonic amplitudes are fixed
and speed-independent, so no single loudness-to-speed mapping holds across
the pool --- and apply the label-fixed augmentations (gain, polarity, channel
drop) from the start of training. Models trained purely on this synthetic
pool, with no real audio, reach 17.8--25.4 PIT-MSE on real free-flight
validation, against $\sim$7.3 for real-data training on its own validation
protocol; a short real-data fine-tune brings them to 11.1--14.1. One honest
note on the history of this result: our first version of this experiment
appeared to fail completely (negative $R^2$), and the failure traced to a
contaminated validation set --- ground-idle segments of one recording labeled
as flight. The trap is generic: at low speeds the evaluation data is as
fragile as the models.

\textbf{The curriculum result.} The transfer result above used only the
label-fixed augmentations. We then combined the strongest generator (refined
annotations on both drones, per-rotor embeddings, per-harmonic checkpoint
selection, maximally randomized conditioning) with the full label-transforming
schedule of Section 6, in a two-stage curriculum: first train on synthetic
noise only (an even mixture of the generator and the analytic comb, full
flight-envelope speed trajectories), then fine-tune on the real data with the
identical schedule. The curriculum improves on its real-only control for
every architecture, under the same frozen validation split and the same
schedule:

\begin{table}[t]
  \centering
  \caption{Generated-first curriculum against real-only training. Frozen
    full-envelope validation split; best-over-epochs, PIT-matched;
    MSE in rev$^2$/s$^2$, MAE in rev/s.}
  \label{tab:curriculum}
  \footnotesize
  \setlength{\tabcolsep}{4pt}
  \begin{tabular}{@{}lccc@{}}
    \toprule
    Architecture & Curriculum & Real-only & $\Delta$MSE \\
    \midrule
    Convolutional & \textbf{28.4} / 3.16 & 52.5 / 3.98 & $-46\%$ \\
    GRU head & 51.4 / 4.14 & 59.2 / 4.24 & $-13\%$ \\
    Transformer head & 38.6 / 3.41 & 42.3 / 3.76 & $-9\%$ \\
    \bottomrule
  \end{tabular}
\end{table}

A per-frame breakdown by flight regime locates the gain: stopped rotors
improve for every architecture (for the convolutional model, silence error
falls from 11.8 to 4.8 rev/s MAE --- nearly its whole aggregate gain), the
warm-up and ramp regimes are mixed across architectures, and mid-flight
changes by at most $\pm 10\%$. The curriculum buys coverage of the regimes
the real corpus lacks, silence above all, at approximately zero cost at
cruise. One honest caveat: for the Transformer, an earlier real-only arm with
a narrower frequency-scaling range reaches 37.6, tying its curriculum row;
the comparison in the table holds the schedule fixed.

\textbf{Ablation: the analytic comb suffices.} Replacing the pre-training
pool by the analytic comb alone (the neural generator removed, everything
else identical) matches or improves every row of the table: 21.1 for the
convolutional model, 34.7 for the Transformer, 51.7 for the GRU. The neural
generator is therefore unnecessary for this downstream gain; the active
ingredients are the guaranteed harmonic structure, the full-envelope speed
trajectories with exact silence, and the label-transforming augmentations.
The pre-training stage's own validation error is uninformative about the
outcome: the Transformer's comb-only pre-training does not transfer at all
before fine-tuning (validation MSE $\sim$2500) and still yields the best
Transformer result after it.

\textbf{Ablation: the staging is necessary.} Merging the real and synthetic
sources into one pool (half real, one quarter generated, one quarter comb;
the same augmentation schedule, one stage instead of two) makes every
architecture worse than its real-only control: 103.8 against 42.3 for the
Transformer, 179.7 against 59.2 for the GRU, and 93.6 against 52.5 for
the convolutional model. Synthetic data
therefore helps only as an initialization that the real stage then refines;
mixed into the real pool it acts as label noise.

\section{Discussion: what the neural models add}
\label{sec:discussion}

Collecting the comparison started in Section 4.2:

\begin{itemize}
\item \textbf{Cruise precision} favors signal processing: 0.69--1.03 rev/s
  blind, against 2.3--2.9 rev/s for the best neural rows (window protocol).
\item \textbf{Coverage} favors the neural models. A ramp-capable variant of
  the two-stage method (coarser at cruise: 1.81/2.52\revps{} against the
  specialist's 0.69/1.03) reaches 3.61\revps{} on the warm-up windows of the
  window protocol, while the neural-initialized rows sit at
  2.24--2.35\revps{} there; on ground and near-silence the tracking method
  emits no estimate at all, while the neural models still answer (with the
  10--15\revps{} silence error of Section 5).
\item \textbf{Latency and cost} favor the neural models: one forward pass per
  frame against per-window iterative optimization (seconds to minutes per
  16-second window on CPU). Applications that need speeds online favor the
  neural path.
\item \textbf{A hybrid} pays only when the initialization is already close.
  On the full window protocol, refining the neural estimates with stage 2
  trims 5--9\%: $2.86 \to 2.68$ and $2.39 \to 2.14$\revps{} for the
  convolutional model, $4.25 \to 4.15$ and $2.26 \to 2.07$ for the GRU
  (DREGON / MD2 cruise). The refinement's capture range bounds the gain: an
  initialization that is already under 1\revps{} (a preliminary two-window
  test with a more accurate recurrent initializer: $0.87 \to 0.64$) is
  refined much further. A sub-rev/s neural initializer is the requirement
  for the hybrid to add value.
\end{itemize}

Limitations: two drone models; tachometer-grade measured speeds exist only
for a subset of DREGON recordings; most numbers come from single training
seeds; all flights are indoor (DREGON) or calm outdoor (MD2); and the
annotation corrections of Section 3.3 are themselves estimated from the
audio, so they share its assumptions --- in particular, every DREGON number
carries the 0.3--0.7 rev/s uncertainty floor.

\section{Conclusion}
\label{sec:conclusion}

Can we train a model to predict motor speeds from drone audio, robustly? On
scarce annotated data, the answer so far has three parts. First, plain
training produces models whose held-out error looks good while their
predictions rely on loudness and the speed prior; the frequency-scaling
probe of Section 6 measures this directly and costs nothing, and we suggest
it as a sanity check for models of this kind. Second, label-transforming
augmentations and full-envelope training data repair part of the problem at
a measurable cost in cruise precision, and a training-free two-stage
tracking method remains 2--4$\times$ more precise at cruise --- a reference we
think learning approaches on this task should report against. Third,
structured synthesis can manufacture the missing diversity only when the
annotations are refined to comb precision and checkpoints are selected by a
per-harmonic measurement; under those conditions, a generated-first
curriculum with the same augmentation schedule improves on real-only
training for every architecture we evaluated, and the improvement sits in
the regimes the real corpus cannot teach --- silence above all.

\appendices

\section{Planned figure list}
\label{app:figures}

\begin{enumerate}
\item Task figure: spectrogram of drone noise with the four-rotor comb and
  the telemetry/refined overlays (Section 3).
\item Frequency-scaling probe: prediction response vs input scaling
  (Section 6).
\item Per-regime error bars: cruise-trained vs full-envelope vs synthetic
  training data (Section 5).
\item Generator line measurement: aggregate loss and line measurement across
  epochs, with spectrograms of the two checkpoint selections (Section 7).
\item Baseline scoreboard: precision vs coverage for all methods (Section 9).
\end{enumerate}

\section{Experiments backing each claim}
\label{app:experiments}

Internal bookkeeping for the authors; strip before submission. Each claim
above maps to an experiment log in \texttt{docs/experiments/}; the mapping
table lives in \texttt{writing/papers/2026-08\_wrapup/inventory.md}.

% Every citation in this draft is still a \pending{cite: ...} marker, so the
% document has no \cite command yet and BibTeX would emit an empty
% bibliography (a LaTeX error). src/references.bib is in place, copied from
% ../../2026-08_joint-decomposition-rps/. Uncomment the two lines below as
% soon as the first real \cite lands.
% \bibliographystyle{IEEEtran}
% \bibliography{references}

\end{document}
